\documentclass[11pt]{article}
\usepackage[T1]{fontenc}
\usepackage{lmodern,microtype,amsmath,amssymb,amsthm,needspace}
\usepackage[margin=1in]{geometry}
\usepackage[colorlinks=true,linkcolor=blue,urlcolor=blue,citecolor=blue]{hyperref}
\hypersetup{pdftitle={A Robinson counterexample to ideal-only tangency certificates},pdfauthor={Alec Kriebel}}
\newtheorem{theorem}{Theorem}
\newtheorem{lemma}[theorem]{Lemma}
\newcommand{\SOS}{\sum S^2}
\title{A Robinson counterexample to\\ideal-only tangency certificates}
\author{Alec Kriebel\\\small Independent researcher\\\small\href{https://orcid.org/0009-0001-9320-500X}{ORCID: 0009-0001-9320-500X}}
\date{12 October 2026 \quad Version 1.0\\\small Unrefereed research note}
\begin{document}
\maketitle
\begin{abstract}
The classical Robinson sextic is nonnegative, but adding any nonnegative
constant does not make it a sum of polynomial squares modulo the ideal of
polynomials vanishing on its real tangency variety. This rules out both the
exact radical certificate and the positive-constant certificate in the proposed
unrestricted replacement of the gradient ideal. The same obstruction applies
to the separate-cross-product ideal printed in the original report. Ten zero
lines force the low-degree homogeneous parts of all square summands to vanish;
an additional coordinate axis then makes the certificate impossible. The proof
allows arbitrarily many summands of arbitrary degree and is entirely elementary.
We credit the classical form and zero-set method, and make no claim of first
bibliographic priority.
\end{abstract}

\section{The proposed certificates}
For $S=\mathbb R[x,y,z]$, let $\SOS$ denote finite sums of squares of elements of
$S$. For $f\in S$, write $f_x,f_y,f_z$ for its partial derivatives and set
\[
 J_f=(yf_x-xf_y,\;zf_x-xf_z,\;zf_y-yf_z),\qquad
 K_f=(yf_x,xf_y,zf_x,xf_z,zf_y,yf_z).
\]
For an ideal $L\subseteq S$, define
\[
 V_{\mathbb R}(L)=\{u\in\mathbb R^3:h(u)=0\text{ for every }h\in L\},\qquad
 I_{\mathbb R}(L)=\{h\in S:h|_{V_{\mathbb R}(L)}=0\}.
\]
The ordinary radical is $\sqrt L=\{h:h^N\in L\text{ for some }N\ge1\}$.
In particular $L\subseteq\sqrt L\subseteq I_{\mathbb R}(L)$.

Schweighofer's Conjecture~2 \cite[p.~811]{report} proposes removing the
minimum-attainment assumption from the gradient-ideal equivalence in
Theorem~1 \cite[p.~810]{report} by replacing that ideal. The four conditions
are global nonnegativity, nonnegativity on the real variety, membership in
$\SOS+\sqrt L$, and membership of $f+\varepsilon$ in $\SOS+L$ for every
$\varepsilon>0$. The printed replacement displays separate products, giving
$K_f$ in three variables. The later catalogue formulation (UnsolvedMath record 30000717)\footnote{\url{https://huggingface.co/datasets/ulamai/UnsolvedMath}.} uses the tangency
minors, giving $J_f$. We treat both explicit readings; we do not assert an
author-approved correction of the printed formula. Removing an attainment
assumption includes polynomials whose minimum \emph{is} attained.

\Needspace{10\baselineskip}
\begin{theorem}\label{thm:main}
Let
\begin{align*}
 R={}&x^6+y^6+z^6-x^4y^2-x^2y^4-x^4z^2-x^2z^4\\
    &\hspace{32mm}-y^4z^2-y^2z^4+3x^2y^2z^2.
\end{align*}
Then $R\ge0$ on $\mathbb R^3$, but for $L=J_R$ or $L=K_R$ and every real
$\varepsilon\ge0$,
\[
 R+\varepsilon\notin\SOS+I_{\mathbb R}(L).
\]
Consequently $R\notin\SOS+\sqrt L$, and
$R+\varepsilon\notin\SOS+L$ for every $\varepsilon>0$.
\end{theorem}
Thus both proposed certificate implications from global nonnegativity fail.
Robinson's form is classical; it already appears in the same report
\cite[p.~798, equation~(2)]{report}. Our purpose is its explicit application
to these ideal-only certificates. The classical non-SOS property alone would
not suffice, since an ideal remainder could change a representation.

\section{The elementary zero-set obstruction}
Put $a=x^2$, $b=y^2$, $c=z^2$ and write $R=F(a,b,c)$. Direct expansion gives
\[
 F=(a-b)^2(a+b-c)+c(a-c)(b-c).
\]
The symmetric polynomial $F$ may be evaluated after ordering $a\ge b\ge c\ge0$;
both displayed terms are then nonnegative. Hence $R\ge0$, and its infimum is
zero, attained at the origin.

Let $Z$ consist of the ten vectors
\[
 (1,s,t),\quad (1,s,0),\quad (1,0,t),\quad (0,1,t),
 \qquad s,t\in\{-1,1\},
\]
where each applicable sign varies independently. Substitution gives $R(v)=0$
for every $v\in Z$, and homogeneity gives $R(uv)=0$ for every real $u$.

\begin{lemma}\label{lem:cubics}
A homogeneous cubic vanishing on $Z$ is zero. The same holds in degrees zero,
one and two.
\end{lemma}
\begin{proof}
The six vectors with a zero coordinate force the cubic to have the form
\[
 p=A x(x^2-y^2-z^2)+B y(y^2-x^2-z^2)
   +C z(z^2-x^2-y^2)+Dxyz.
\]
At $(1,s,t)$ its value is $-A-Bs-Ct+Dst$. Summing these four values against
the sign characters $1,s,t,st$ gives $-4A,-4B,-4C,4D$, respectively. Thus all
four coefficients vanish. If a homogeneous degree-$d$ polynomial $p_d$, with
$0\le d<3$, vanishes on $Z$, then $x^{3-d}p_d$ is such a cubic. Since $S$ is an
integral domain, $p_d=0$.
\end{proof}
The ten-zero cubic obstruction is classical; Reznick gives the same ten
directions and cubic argument in \cite[Section~4(b), p.~259]{reznick}.
We include the full argument, then use whole
zero lines and an additional axis to address the ideal remainder and positive
constant perturbations.

Every point of the lines $\mathbb Rv$, $v\in Z$, is a global minimizer of $R$.
Its gradient is therefore zero, and each line lies in both $V_{\mathbb R}(J_R)$
and $V_{\mathbb R}(K_R)$. Both varieties also contain the axis $(u,0,0)$:
there $\nabla R=(6u^5,0,0)$, all cross-products vanish, and $R=u^6$.

\begin{proof}[Proof of Theorem~\ref{thm:main}]
Fix $\varepsilon\ge0$ and either ideal $L$. Suppose
\[
 R+\varepsilon=\sum_{j=1}^m q_j^2+h,
 \qquad q_j\in S,\quad h\in I_{\mathbb R}(L).
\]
On every zero line, $h$ vanishes and
$\varepsilon=\sum_j q_j(uv)^2$ for all real $u$. Hence each real univariate
polynomial $q_j(uv)$ is bounded in absolute value by $\sqrt\varepsilon$ on
$\mathbb R$, and must be constant. Write $q_j=\sum_d q_{j,d}$ in homogeneous
parts. Comparing powers of $u$ gives $q_{j,d}(v)=0$ for every $d\ge1$ and
$v\in Z$. Lemma~\ref{lem:cubics} yields
$q_{j,1}=q_{j,2}=q_{j,3}=0$.

On the coordinate axis the certificate gives
\[
 u^6+\varepsilon=\sum_j q_j(u,0,0)^2.
\]
Every nonzero restriction on the right has degree at most three. Indeed, if
$D>3$ is the largest such degree, the coefficient of $u^{2D}$ is a nonempty
sum of squares of nonzero real leading coefficients, which cannot cancel.
But the degree-one, degree-two and degree-three homogeneous parts have all
vanished identically. The restrictions must therefore be constant, contradicting
the nonconstant left side. The inclusions preceding the theorem give its
remaining conclusions.
\end{proof}
Constants were retained throughout: when $\varepsilon>0$ the $q_j$ need not
vanish on the zero lines. The proof imposes no global degree bound on the
summands, radicality hypothesis, or finiteness condition on the varieties.

\section{Formulation boundaries}
The following standard sphere-minimization argument also underlies the
tangency-set infimum identity in \cite[Section~7.2.3, Lemma~7.21]{laurent}.
For any polynomial $f$ in $n$ variables, nonnegativity on all of $\mathbb R^n$
is equivalent to nonnegativity on the real variety of its tangency minors.
If $f(w)<0$ and $w=0$, the origin already belongs to that variety. Otherwise
minimize $f$ on the sphere of radius $\|w\|$. At a minimizer $v$, Lagrange
multipliers give $\nabla f(v)=\lambda v$ and $f(v)\le f(w)<0$, so all tangency
minors vanish. This proves the reverse implication without any assumption
about the global infimum. For the Robinson example, the first two proposed
conditions therefore hold and both certificates fail.

The literal product reading has a further, separate defect. We use a constant
shift of the standard nonattaining example in
\cite[Section~1.6, equation~(6)]{tentacles} and \cite[p.~810]{report}:
\[
 Q(x,y)=(1-xy)^2+y^2-\tfrac12,\qquad K_Q=(yQ_x,xQ_y).
\]
Here $yQ_x=2y^2(xy-1)$ and $xQ_y=2x(x(xy-1)+y)$. The first equation forces
$y=0$ or $xy=1$. The second then forces $x=0$ in the former case, and equals
$2$, a contradiction, in the latter. Thus $V_{\mathbb R}(K_Q)=\{(0,0)\}$.
Nevertheless $Q(0,0)=1/2$ and $Q(2,1/2)=-1/4$. Its infimum $-1/2$ is not
attained: equality would require both $y=0$ and $xy=1$, whereas
$Q(1/u,u)=u^2-1/2$ approaches the infimum. This diagnostic applies only to the
product reading; it does not negate the tangency-minor equivalence just proved.

For the actual gradient ideal $G_R=(R_x,R_y,R_z)$, Euler's identity gives
$6R=xR_x+yR_y+zR_z$, hence $R\in G_R$. There is no contradiction with the
gradient-ideal theorem. Likewise the gradient-tentacle inequalities
\cite{tentacles} and truncated tangency certificates
\cite{truncated,gcv} use additional inequalities and SOS-multiplier terms;
they are different certificates, which this note does not refute.

\section*{Verification, attribution and review status}
The accompanying source archive preserves the original candidate and its
historical audits, fresh independent reviews, and exact reproduction controls.
The cubic evaluation determinant is $-128$ in the stated monomial order in
the original proof. Symbolic checks supplement the universal proof; no numerical
SDP failure or fixed SOS-degree cutoff is used as evidence of nonexistence.
Robinson's form and the classical zero-set method retain their prior attribution.
A bounded priority audit is provided separately; first priority and continued
open status through 2026 are not established or asserted. Complete 2008 and
2009 H\`a--Pham article texts were not obtained; comparisons with truncated
certificates rely on later primary statements, and the access gaps remain in
the audit.

AI tools were used extensively in developing, checking and presenting this
argument, including independent AI-assisted adversarial reviews. This note is
unrefereed and has not received independent human or journal peer review.
Alec Kriebel is the author of the submitted note. The authored note and
verification materials are released under CC BY 4.0; cited source documents
retain their own terms and are not redistributed.

\begin{thebibliography}{9}
\bibitem{report}
M.~Schweighofer, \emph{Gradient tentacles and sums of squares}, pp.~810ff.;
B.~Reznick, \emph{Hilbert's construction of psd polynomials that are not a sum
of squares}, pp.~797--799; in \emph{Reelle Algebraische Geometrie},
Oberwolfach Reports 14/2007, pp.~753--820.
\href{https://doi.org/10.4171/owr/2007/14}{doi:10.4171/owr/2007/14}.
\bibitem{reznick}
B.~Reznick, \emph{Some concrete aspects of Hilbert's 17th problem},
Contemporary Mathematics 253 (2000), 251--272.
\href{https://doi.org/10.1090/conm/253/03936}{doi:10.1090/conm/253/03936}.
\href{https://reznick.web.illinois.edu/Paper34.pdf}{Author-hosted published reprint}.
\bibitem{tentacles}
M.~Schweighofer, \emph{Global optimization of polynomials using gradient
tentacles and sums of squares}, SIAM J. Optim. 17 (2006), 920--942.
\href{https://arxiv.org/abs/math/0511617}{arXiv:math/0511617}.
\bibitem{truncated}
H\`a Huy Vui and Pham Tien Son, \emph{Global optimization of polynomials using
the truncated tangency variety and sums of squares}, SIAM J. Optim. 19 (2008),
941--951. \href{https://doi.org/10.1137/080719212}{doi:10.1137/080719212}.
\bibitem{laurent}
M.~Laurent, \emph{Sums of squares, moment matrices and optimization over
polynomials}, in M.~Putinar and S.~Sullivant (eds.), \emph{Emerging Applications
of Algebraic Geometry}, IMA Volumes in Mathematics and its Applications 149,
Springer, 2009, pp.~157--270.
\href{https://doi.org/10.1007/978-0-387-09686-5_7}{doi:10.1007/978-0-387-09686-5\_7}.
\bibitem{gcv}
F.~Guo, M.~Safey El Din and L.~Zhi, \emph{Global optimization of polynomials
using generalized critical values and sums of squares}, ISSAC 2010.
\href{https://perso.lip6.fr/Mohab.Safey/Articles/gcv_sos.pdf}{Author manuscript}.
\end{thebibliography}
\end{document}
